\documentclass[../../../thesis.tex]{subfiles}
\begin{document}
\subsection{Odkazy na citované diela}
Vo vedecko-technických oblastiach, do ktorých patria aj študijné programy na našej fakulte, je zvykom používať numerický systém citovania. Jednotlivé zdroje sú očíslované podľa poradia výskytu v texte, pričom číslo zdroja uvádzame v hranatých zátvorkách.
Citovanie sa riadi technickou normou STN ISO 690: 2022 Informácie a dokumentácia: Návod na tvorbu bibliografických odkazov na informačné pramene a~ich citovanie~\cite{iso690}.

Systém \LaTeX\ umožňuje pracovať s citáciami niekoľkými spôsobmi.
Šablóna FEIstyle pracuje s balíčkom \verb|biblatex| na správu citácií. Jeho základom je externý databázový súbor \texttt{bibliography.bib},
ktorý sa nachádza v hlavnom priečinku projektu záverečnej práce.
Súbor obsahuje bibliografické záznamy citovaných diel
v špecifickom formáte.
Každý záznam začína jedinečným identifikátorom, ktorý zvyčajne
volíme tak, aby sme si ho jednoducho pamätali.
Ak totiž v texte chceme dielo citovať, napíšeme makro \verb|\cite{identifikator}| a~v~práci sa automaticky objaví poradové číslo citovaného diela v hranatej zátvorke.
Číslo citácii pridelí externý program \textsl{biber},
ktorý prejde celý dokument, identifikuje výskyty makra \verb|\cite|, zoradí citované diela podľa výskytu v~práci,
pridelí im čísla a vytvorí podklady na zoznam literatúry na konci práce.
Pri ďalšom spustení kompilátora \TeX-u dôjde k~nahradeniu makier
\verb|\cite| číslami v hranatých zátvorkách a makro šablóny
\verb|\FEIbibliography| vytvorí zoznam literatúры s nadpisom
prvej úrovne Literatúra.

Pri použití bibliografie odporúčame použiť dávkový súbor
\texttt{Makefile} alebo spustiť externý progрam \texttt{biber}
po prvej kompilácii a~potom skompilovať dokument ešte dvakrát.
Posledná kompilácia zabezpečí správne čísla strán v obsahu.
Online nástroj Overleaf robí všetky potrebné kroky v jednom behu,
kompiláciu tu stačí spustiť iba raz.
\end{document}